%% JAIR Manuscript (acmart) - Clement Frassier
%% Federated Learning for Stress Detection on WESAD (wrist-only, TinyML)
%% Complementary case study to the nurse-stress-detection report

\documentclass{acmart}
\usepackage{amsmath, amssymb}
\usepackage{booktabs}
\usepackage{graphicx}
\usepackage{hyperref}
\usepackage{xcolor}

\RequirePackage[
  datamodel=acmdatamodel,
  style=numeric,
  backend=biber,
  giveninits=true,
  uniquename=init
  ]{biblatex}

\renewcommand*{\bibopenbracket}{(}
\renewcommand*{\bibclosebracket}{)}

\addbibresource{sample-base.bib}

\begin{document}

\title[FL sur WESAD : cas de contrôle IID]{Apprentissage fédéré pour la détection du stress sur WESAD : un cas de contrôle quasi-IID face à l'hétérogénéité clinique réelle}

\author{Clement Frassier}
\orcid{0009-0000-0000-0000}
\email{clement.frassier@gmail.com}
\affiliation{%
  \institution{UKM / UEL}
  \city{Kuala Lumpur}
  \country{Malaysia}
}

\begin{abstract}
{\bf Background:}
Le rapport principal de ce projet (détection du stress chez 15 infirmières
hospitalières) a montré qu'une hétérogénéité clinique extrême limite la
généralisation de toute architecture, centralisée comprise. Ce document teste
la même question méthodologique sur un second jeu de données, WESAD
\cite{schmidt2018wesad}, collecté en laboratoire selon un protocole contrôlé et
identique pour tous les sujets, pour vérifier si le coût mesuré du FL persiste
quand l'hétérogénéité entre clients est bien plus faible.

{\bf Objectives:}
Reproduire, sur WESAD, le même pipeline d'apprentissage fédéré (FedProx +
DP-SGD, puis FedPer) que celui validé sur le jeu de données infirmières, en se
restreignant aux capteurs du poignet (Empatica E4) pour rester dans une
contrainte de déploiement TinyML réaliste, et comparer les résultats à une
baseline centralisée non contrainte.

{\bf Methods:}
15 sujets, capteurs poignet uniquement (BVP 64~Hz, EDA 4~Hz, température 4~Hz,
accélération 32~Hz, tous ramenés à 4~Hz commun), fenêtrage glissant identique
au pipeline infirmières (60/30 échantillons, statistiques
moyenne/écart-type/min/max), label binarisé stress (protocole TSST) contre
non-stress (baseline, amusement, méditation). Même architecture MLP Tiny
(1~362 paramètres, 5,32~Ko). Grille complète $5\sigma \times 5$ graines pour
FedProx+DP et pour FedPer (50 runs fédérés), plus 5 graines centralisées.

{\bf Results:}
Sur les 8 clients au test-set réellement mixte (7 des 15 sont mono-classe par
construction chronologique du protocole, pas par hétérogénéité clinique),
\textbf{aucun effondrement n'est observé} : balanced-accuracy
$79{,}6 \pm 6{,}7\,\%$ pour FedProx+DP, $78{,}3 \pm 1{,}7\,\%$ pour le
centralisé (écart non significatif, $p=0{,}44$), $87{,}6 \pm 3{,}8\,\%$ pour
FedPer (significativement supérieur aux deux autres, $p<10^{-5}$). L'AUC-ROC
reste toujours nettement au-dessus du hasard ($0{,}88$ à $0{,}91$ selon la
condition), à l'opposé du $0{,}267$ mesuré sur les clients difficiles du jeu de
données infirmières.

{\bf Conclusions:}
Sur un protocole de laboratoire standardisé, quasi-IID entre sujets, le FL
n'ajoute aucun coût statistiquement mesurable par rapport au centralisé, et la
personnalisation (FedPer) fait mieux que les deux architectures non
personnalisées. Ce résultat renforce, par contraste, la lecture du rapport
principal : ce n'est pas la fédération en tant que telle qui limite la
généralisation sur le jeu de données infirmières, mais l'hétérogénéité
clinique réelle propre à ce jeu de données, absente ici par construction du
protocole.
\end{abstract}

\received{6 July 2026}

\maketitle

\section{Introduction}

Le rapport principal de ce projet établit un résultat central sur le jeu de
données infirmières : l'hétérogénéité clinique extrême entre les 15
participantes limite la généralisation de toute architecture testée,
centralisée y compris, et le FL n'ajoute qu'un désavantage marginal par-dessus
cette limite de fond. Une question naturelle en découle : ce coût mesuré,
faible mais réel, tient-il à la fédération elle-même ou à l'hétérogénéité
propre à ce jeu de données précis ?

Ce document répond en rejouant le même pipeline sur \textbf{WESAD}
(\textit{Wearable Stress and Affect Detection}) \cite{schmidt2018wesad}, un
jeu de données de 15 sujets enregistrés en laboratoire selon un protocole
identique et contrôlé (baseline, test de stress TSST, amusement, deux phases
de méditation). Contrairement au recrutement clinique réel des infirmières,
WESAD est \textbf{quasi-IID entre sujets} par construction : la composition
de classe (stress~$\approx 22\,\%$) varie de moins de trois points d'un sujet
à l'autre. C'est précisément ce contraste qui rend la comparaison utile : si
le coût du FL mesuré sur les infirmières était un artefact générique de la
fédération, on devrait le retrouver ici ; s'il ne l'est pas, cela confirme
que l'hétérogénéité clinique, pas le FL, était la variable en cause.

Ce travail reste volontairement plus resserré que le rapport principal (une
seule architecture de personnalisation testée, pas d'ablation FedRep+SCAFFOLD
ni de scénario nouvel arrivant) : son rôle est celui d'un \textbf{cas de
contrôle}, pas d'une étude indépendante de même ampleur.

\section{Méthodologie}
\label{sec:method}

\subsection{Jeu de données et restriction aux capteurs du poignet}

WESAD~\cite{schmidt2018wesad} enregistre 15 sujets (S2--S17, S1 et S12
exclus par les auteurs pour défaillance capteur) via deux dispositifs : un
capteur torse RespiBAN (ECG, EMG, respiration, EDA, température,
accélération, 700~Hz) et un bracelet Empatica~E4 au poignet (BVP 64~Hz, EDA
4~Hz, température 4~Hz, accélération 32~Hz). Le RespiBAN est un appareil de
recherche, pas un dispositif portable réaliste ; pour rester cohérent avec la
contrainte TinyML du rapport principal et permettre une comparaison directe,
\textbf{seuls les signaux du poignet sont retenus}.

Les quatre canaux du poignet, à des fréquences d'échantillonnage différentes,
sont ramenés à une fréquence commune de 4~Hz (celle de l'EDA et de la
température) par moyennage par blocs non chevauchants (facteur~8 pour
l'accélération à 32~Hz, facteur~16 pour le BVP à 64~Hz). Le label de protocole
(échantillonné à 700~Hz sur le flux torse, synchronisé en $t=0$ avec le
poignet) est ramené à la même fréquence par vote majoritaire par bloc de 175
échantillons.

Le label brut comporte 8 valeurs : 0~=~transitoire, 1~=~baseline,
2~=~stress, 3~=~amusement, 4~=~méditation, 5/6/7~=~à ignorer (voir la
documentation du jeu de données). Seules les valeurs 1 à 4 sont conservées
(comme le fait le protocole officiel), puis binarisées : stress~=~1 (label
brut~=~2), non-stress~=~0 (labels bruts~1, 3, 4). Ce filtrage a lieu
\emph{avant} le fenêtrage, ce qui peut introduire des discontinuités
temporelles à la frontière des segments écartés, exactement comme dans le
pipeline infirmières.

Le fenêtrage glissant est identique à celui du rapport principal : fenêtre de
60 échantillons (15~s à 4~Hz), pas de 30 (50~\% de recouvrement), 4
statistiques (moyenne, écart-type, min, max) par canal, sur les 6 canaux
(accélération~x/y/z, BVP, EDA, température) — soit \textbf{24 features par
fenêtre}, la même dimensionnalité que le pipeline infirmières. Au total,
\textbf{5~971 fenêtres} sur les 15 sujets ($\approx 398$ par sujet), un
ordre de grandeur bien plus petit que les 383~584 fenêtres du jeu de données
infirmières.

\paragraph{Split temporel et une mono-classe d'origine différente.} Le split
train/test (80/20, gap de 2 fenêtres) est chronologique par sujet, comme dans
le rapport principal. Mais ici, \textbf{7 des 15 sujets ont un test-set
mono-classe} (uniquement non-stress), \emph{non} par hétérogénéité clinique
entre individus (la composition de classe globale est quasi identique d'un
sujet à l'autre, Section~\ref{sec:results}), mais parce que le protocole WESAD
se termine systématiquement par les phases de méditation
(non-stress) : pour de nombreux sujets, la tranche chronologique finale
(20~\% des fenêtres utilisables) tombe entièrement dans cette phase de fin de
protocole. C'est un artefact de \textbf{l'ordre du protocole}, pas de
\textbf{l'hétérogénéité entre sujets} — une distinction importante détaillée
en Section~\ref{sec:discussion}. Les 8 sujets au test-set mixte
(indices 2, 3, 5, 6, 8, 10, 12, 14) sont les seuls porteurs d'un signal
informatif sur la généralisation du modèle, et toute métrique groupée sur les
15 sujets est rapportée à côté de la métrique restreinte à ces 8 sujets, selon
la même discipline que le rapport principal.

\subsection{Modèle, FedProx, DP-SGD et FedPer}

Le modèle (MLP Tiny, 1~362 paramètres, 5,32~Ko), la stratégie FedProx
($\mu=0{,}05$), le mécanisme DP-SGD (Opacus~\cite{opacus}, clipping
$C=1{,}0$, $\delta=10^{-5}$) et l'architecture de personnalisation FedPer
(extracteur \texttt{fc1}/\texttt{fc2} agrégé, tête \texttt{fc3} jamais
transmise) sont \textbf{strictement identiques} au rapport principal (voir
\texttt{rapport\_acc\_nouveau\_.tex}, Sections~2.3 à 2.6, pour le détail des
formules). Seul le prétraitement des données change.

\subsection{Protocole expérimental}

Simulation Flower~\cite{beutel2020flower}, 15 clients. Grille complète
$\sigma \in \{0{,}6;\,0{,}8;\,1{,}0;\,1{,}4;\,1{,}8\}$ $\times$ 5 graines (42,
123, 7, 2024, 31415), 30 rounds, 2 époques locales par round — pour
\textbf{FedProx+DP} et pour \textbf{FedPer} séparément (50 runs fédérés au
total), plus une baseline centralisée sur les mêmes 5 graines (standardisation
globale, pas de FL, pas de DP). L'évaluation FedPer utilise la méthodologie
corrigée du rapport principal : l'extracteur partagé du checkpoint global final
est combiné à la tête locale propre à chaque client (pas la copie dérivée de
l'extracteur que chaque client a localement), pour éviter le biais déjà
identifié et corrigé dans le rapport principal.

\section{Résultats}
\label{sec:results}

\subsection{Composition de classe par sujet}

\begin{table}[h!]
  \centering
  \small
  \caption{Composition du test-set par sujet (proportion de stress). 7/15
  sujets sont mono-classe (0~\% stress), par construction chronologique du
  protocole (Section~\ref{sec:method}), pas par hétérogénéité clinique.}
  \label{tab:class_composition}
  \begin{tabular}{cccc}
    \toprule
    Sujet & \% stress (test) & Sujet & \% stress (test) \\
    \midrule
    0 (S2)  & 0,0\% (mono-classe)  & 8 (S10)  & 35,0\% \\
    1 (S3)  & 0,0\% (mono-classe)  & 9 (S11)  & 0,0\% (mono-classe) \\
    2 (S4)  & 32,5\%               & 10 (S13) & 34,2\% \\
    3 (S5)  & 33,3\%               & 11 (S14) & 0,0\% (mono-classe) \\
    4 (S6)  & 0,0\% (mono-classe)  & 12 (S15) & 34,2\% \\
    5 (S7)  & 33,3\%               & 13 (S16) & 0,0\% (mono-classe) \\
    6 (S8)  & 33,3\%               & 14 (S17) & 34,6\% \\
    7 (S9)  & 0,0\% (mono-classe)  & &  \\
    \bottomrule
  \end{tabular}
\end{table}

Chez les 8 sujets au test-set mixte, la proportion de stress varie de
$32{,}5\,\%$ à $34{,}6\,\%$ : un écart de moins de 3 points, à comparer aux
dizaines de points d'écart mesurés entre clients du jeu de données infirmières.
C'est la confirmation empirique que WESAD est bien un cas quasi-IID.

\subsection{FedProx+DP, FedPer et centralisé : pas d'effondrement}

\begin{table}[h!]
  \centering
  \small
  \caption{Comparaison des trois conditions, métrique groupée sur les 15
  sujets contre restreinte aux 8 sujets au test-set mixte (moyenne $\pm$
  écart-type sur 25 runs FL [$5\sigma\times$5 graines] ou 5 graines
  centralisées).}
  \label{tab:main_wesad}
  \begin{tabular}{lccc}
    \toprule
    Balanced-accuracy & FedProx+DP & FedPer & Centralisé \\
    \midrule
    Groupée (15 sujets)    & $87{,}67 \pm 2{,}91\,\%$ & $\mathbf{90{,}72 \pm 1{,}70\,\%}$ & $82{,}87 \pm 0{,}89\,\%$ \\
    \textbf{Sujets mixtes (8/15)} & $79{,}61 \pm 6{,}72\,\%$ & $\mathbf{87{,}63 \pm 3{,}83\,\%}$ & $78{,}33 \pm 1{,}71\,\%$ \\
    \midrule
    AUC-ROC (sujets mixtes) & $0{,}8897 \pm 0{,}0055$ & $\mathbf{0{,}9085 \pm 0{,}0184}$ & $0{,}8843 \pm 0{,}0257$ \\
    Accuracy (sujets mixtes) & $76{,}79 \pm 8{,}91\,\%$ & $\mathbf{87{,}89 \pm 5{,}08\,\%}$ & $78{,}37 \pm 2{,}61\,\%$ \\
    \bottomrule
  \end{tabular}
\end{table}

\textbf{Aucune des trois conditions ne s'effondre} sur les sujets mixtes :
l'AUC-ROC reste dans tous les cas entre $0{,}88$ et $0{,}91$, loin du
$0{,}267$ mesuré sur les clients difficiles du jeu de données infirmières. La
balanced-accuracy groupée sur les 15 sujets est gonflée par les 7 sujets
mono-classe dans les trois conditions (de 6 à 8 points selon la condition),
le même phénomène de distorsion déjà documenté dans le rapport principal, mais
sans jamais renverser le signe de la conclusion : les modèles généralisent
réellement sur ce jeu de données, contrairement au cas infirmières.

\paragraph{Tests statistiques (sujets mixtes, Welch et Cohen's~$d$).}
\begin{itemize}
  \item \textbf{Centralisé vs FedProx+DP : pas de différence significative}
  (balanced-accuracy $t=-0{,}79$, $p=0{,}44$, $d=-0{,}25$ ; AUC $t=-0{,}42$,
  $p=0{,}70$, $d=-0{,}26$). Contrairement au jeu de données infirmières, la
  fédération et la confidentialité différentielle n'ajoutent ici
  \textbf{aucun coût statistiquement détectable}.
  \item \textbf{FedPer vs FedProx+DP : différence significative et large}
  (balanced-accuracy $t=5{,}08$, $p=1{,}0\times10^{-5}$, $d=1{,}44$ ; AUC
  $t=4{,}79$, $p=4{,}8\times10^{-5}$, $d=1{,}36$). La personnalisation aide
  nettement, comme dans le rapport principal.
  \item \textbf{FedPer vs Centralisé : différence significative sur la
  balanced-accuracy} ($t=8{,}01$, $p=3{,}6\times10^{-6}$, $d=3{,}02$, effet
  très large), \textbf{mais pas sur l'AUC} ($t=1{,}81$, $p=0{,}13$,
  $d=1{,}00$) — l'effectif réduit du centralisé ($n=5$) limite la puissance
  statistique sur cette métrique précise ; le signe de l'effet reste
  cohérent dans les deux cas (FedPer au-dessus).
\end{itemize}

\paragraph{Le même piège de la métrique groupée, dans l'autre sens.} Sur la
métrique \emph{groupée} (15 sujets), Centralisé vs FedProx+DP est
significatif ($t=-6{,}48$, $p=2{,}4\times10^{-6}$), mais dans le sens opposé
à ce que l'intuition attendrait : FedProx+DP y paraît \emph{meilleur} que le
centralisé. Restreint aux 8 sujets mixtes, cet écart disparaît complètement
($p=0{,}44$). C'est le même mécanisme de distorsion déjà identifié dans le
rapport principal (les sujets mono-classe dominent la métrique groupée), ici
manifesté dans la direction opposée : une confirmation supplémentaire qu'il
ne faut jamais lire une métrique groupée sur des clients mono-classe sans la
croiser avec la métrique restreinte aux clients informatifs, quel que soit le
sens apparent de l'écart.

\subsection{Budget de confidentialité}

\begin{table}[h!]
  \centering
  \small
  \caption{Grille confidentialité/utilité, FedProx+DP (moyenne $\pm$
  écart-type sur 5 graines, round 30, $\delta=10^{-5}$).}
  \label{tab:eps_wesad}
  \begin{tabular}{cc}
    \toprule
    $\sigma$ & $\varepsilon$ (round 30) \\
    \midrule
    0,6 & $35{,}11 \pm 1{,}14$ \\
    0,8 & $18{,}83 \pm 0{,}73$ \\
    1,0 & $13{,}12 \pm 0{,}77$ \\
    1,4 & $7{,}18 \pm 0{,}19$ \\
    1,8 & $5{,}11 \pm 0{,}23$ \\
    \bottomrule
  \end{tabular}
\end{table}

À $\sigma$ égal, $\varepsilon$ est ici \textbf{nettement plus élevé} que sur le
jeu de données infirmières (par exemple $5{,}11$ contre $0{,}54$ à
$\sigma=1{,}8$). Ceci ne traduit pas une confidentialité moins bien
implémentée, mais la taille bien plus petite du jeu d'entraînement local par
client ($\approx 320$ fenêtres contre $\approx 20\,000$ chez les infirmières) :
à nombre d'époques et taille de batch égaux, le ratio d'échantillonnage par
batch est mécaniquement plus élevé, ce qui accroît le budget de
confidentialité consommé par round pour un même $\sigma$. Un $\sigma$ plus
élevé serait nécessaire ici pour viser le même $\varepsilon$ cible que le
rapport principal.

\subsection{Coût de communication}

Le modèle FedProx+DP transmet 5,32~Ko par round (identique au rapport
principal, même architecture). FedPer transmet 5,19~Ko (la tête
\texttt{fc3}, 34 paramètres, n'étant jamais transmise), un gain de
communication de $-2{,}5\,\%$, identique en proportion à celui mesuré sur le
jeu de données infirmières.

\section{Discussion}
\label{sec:discussion}

\subsection{Ce que confirme ce cas de contrôle}

Le rapport principal concluait que l'hétérogénéité clinique extrême entre
infirmières, pas la fédération elle-même, expliquait l'essentiel de l'écart
mesuré avec le centralisé. Ce cas de contrôle sur WESAD apporte une preuve
empirique directe de cette lecture : sur un protocole de laboratoire
standardisé et quasi-IID, \textbf{l'écart FL-vs-centralisé disparaît
statistiquement} (Section~\ref{sec:results}), alors que le pipeline, le
modèle et les mécanismes de confidentialité sont rigoureusement identiques.
Si le coût mesuré sur les infirmières provenait d'une limite générique du FL
ou de DP-SGD, on s'attendrait à retrouver un coût comparable ici ; ce n'est
pas le cas.

\subsection{FedPer dépasse le centralisé : une différence de mécanisme, pas
seulement de degré}

Sur le jeu de données infirmières, la personnalisation (FedPer,
FedRep+SCAFFOLD) réduisait l'écart avec le centralisé sans jamais le faire
disparaître ni le dépasser. Ici, \textbf{FedPer dépasse significativement le
centralisé} sur la balanced-accuracy ($87{,}63\,\%$ contre $78{,}33\,\%$,
$p=3{,}6\times10^{-6}$). Une explication plausible : même sur un protocole
standardisé, chaque individu a sa propre ligne de base physiologique (BVP,
EDA, température au repos varient d'une personne à l'autre) ; un modèle
centralisé unique doit trouver un compromis global, alors qu'une tête de
classification propre à chaque sujet peut absorber cette variabilité
individuelle sans que l'extracteur partagé n'ait besoin de la capturer. Ce
mécanisme de personnalisation profite ici pleinement de la petite taille du
jeu de données locale (peu de risque de sous-apprentissage sur 34 paramètres
ajustés par sujet), un résultat cohérent avec le constat déjà fait dans le
rapport principal que la personnalisation \emph{de la tête seule} est le
facteur dominant, pas la durée d'entraînement.

\subsection{La mono-classe d'origine chronologique n'est pas la mono-classe
d'origine clinique}

Il est important de ne pas confondre les deux phénomènes de mono-classe
rencontrés dans ce projet. Chez les infirmières, la mono-classe résultait
d'une \textbf{hétérogénéité réelle entre individus} (certaines infirmières
n'ont simplement jamais connu d'épisode de stress ou de repos pendant leur
session). Ici, elle résulte d'un \textbf{artefact de l'ordre du protocole}
(le déroulé expérimental se termine systématiquement par des phases de
non-stress) : rien n'indique que ces 7 sujets soient physiologiquement
atypiques, seule la position chronologique de leur test-set dans le protocole
diffère. La distinction compte pour l'interprétation : sur WESAD, un split
non chronologique (par exemple stratifié) aurait probablement éliminé ce
phénomène presque entièrement, ce qui n'aurait pas été le cas pour les
infirmières où l'hétérogénéité est structurelle au jeu de données, pas un
artefact du protocole de split.

\section{Limites et travaux futurs}
\label{sec:limits}

\begin{itemize}
  \item \textbf{Portée de ce cas de contrôle.} Ce document teste une seule
  architecture de personnalisation (FedPer) sur un seul jeu de données de
  contrôle. Il ne prétend pas à l'exhaustivité du rapport principal (pas
  d'ablation FedRep+SCAFFOLD, pas de scénario nouvel arrivant, pas de
  décomposition fédération-vs-confidentialité par une grille FedProx sans
  DP-SGD dédiée).
  \item \textbf{Quantification INT8 non mesurée.} Contrairement au rapport
  principal, cette étude ne mesure pas l'erreur de quantification dynamique
  INT8 ; seule la précision FP32 est rapportée ici.
  \item \textbf{Split chronologique et mono-classe protocolaire.} Comme
  discuté en Section~\ref{sec:discussion}, 7 des 15 sujets ont un test-set
  mono-classe pour une raison propre au protocole WESAD, pas au jeu de
  données infirmières. Un split alternatif (stratifié plutôt que
  chronologique) permettrait de vérifier que les conclusions ne dépendent
  pas de ce choix de split, mais introduirait alors un risque de fuite entre
  fenêtres chevauchantes (Section~\ref{sec:method}) qu'il faudrait traiter
  spécifiquement.
  \item \textbf{Restriction aux capteurs du poignet.} Le RespiBAN (torse) a
  été volontairement écarté pour la contrainte TinyML ; les résultats ne
  sont donc pas directement comparables aux publications WESAD originales
  qui exploitent l'ECG et la respiration du torse.
  \item \textbf{Effectif centralisé réduit ($n=5$).} Comme dans le rapport
  principal, certaines comparaisons impliquant le centralisé (5 graines
  seulement) ont une puissance statistique plus faible que les comparaisons
  FL-vs-FL (25 valeurs) ; le cas de l'AUC FedPer-vs-Centralisé
  (Section~\ref{sec:results}) illustre cette limite.
  \item \textbf{Grille de confidentialité non recalibrée pour ce jeu de
  données.} Le budget $\varepsilon$ obtenu ici à $\sigma$ égal est
  nettement plus élevé que sur le jeu de données infirmières
  (Section~\ref{sec:results}), du fait de la taille bien plus réduite du
  jeu d'entraînement local. Une grille dédiée avec des valeurs de $\sigma$
  plus élevées serait nécessaire pour viser un $\varepsilon$ cible
  comparable au rapport principal.
\end{itemize}

\printbibliography

\end{document}
